\documentclass[11pt,compress,t,notes=noshow, xcolor=table]{beamer}

\input{../../style/preamble}
\input{../../latex-math/basic-math}
\input{../../latex-math/basic-ml}

% local helpers for this chunk
\newcommand{\zv}{\bm{z}}                 % node value (vector)
\newcommand{\adj}[1]{\bar{\bm{z}}_{#1}}  % adjoint of node i

\title{Optimization in Machine Learning}

\begin{document}

\titlemeta{
Advanced first order methods
}{
Backpropagation and autodiff
}{
figure/nn_architecture.pdf
}{
\item DNNs as computational graphs
\item Gradients via the chain rule and vector-Jacobian products (VJPs)
\item Backpropagation algorithm
\item Implementation via automatic differentiation
\vspace{0.5cm}
\item[] \hspace{-0.7cm}Slides based on\\
\hspace{-0.7cm}\furtherreading{JAX2024AUTODIFF} \furtherreading{DAGREOU2024HVP}
}


%TODO-BB: auf jeden fall referenzen notieren

%TODO-BB: vielleicht minimal anderen kontext zum einstieg
%TODO-BB: wir können hier gern über DL reden die ganze zeit, wir sollten aber EINMAL
% klar sagen dass man den autodiff krams halt für mehr als "predictive DL models" nutzen kann....

% vielleich ist es doch besser über beliebige algebraische ausdrücke hier zu reden....

\begin{framei}{Computing gradients for DNNs}
\item We now train \textbf{deep neural networks (DNNs)}: models with a very large parameter vector $\thetav \in \R^{d}$ ($d$ potentially in the millions)
\item All methods discussed in this chapter rely on the gradient of the (scalar) empirical risk w.r.t.\ $\thetav$:
$$
\nabla_{\thetav} \risket \in \R^{1 \times d}
$$
\\ {\footnotesize(SGD relies on mini-batch estimates of the gradient)}
% \item The objective is \textbf{scalar}, so this is a single row gradient -- but $p$ is huge

%TODO-BB: "L" velleicht nicht optimal hier?
\item For brevity, we write $L := \risket$ in the following
% \item The objective is \textbf{scalar}, so this is a single row gradient -- but $p$ is huge
\item For brevity, we write $L := \risket$ in the following
% \item The objective is \textbf{scalar}, so this is a single row gradient -- but $d$ is huge
\spacer
\item \textbf{Question:} How to compute it \textbf{efficiently}?
\end{framei}


%TODO-BB: wir wollen hier wohl theta = (theta_1, .... theta_k) schreiben.
% jedes theta_i ist dann der vektor der params des i-ten layers. sollten wir sagen.
% theta_(i) wäre vielleicht auch klarer.....
% ausserdem noch sagen dass p_i halt die länge von theta_i ist....

\begin{framev}{DNNs as computational graphs}
\begin{itemize}
\item A DNN can be represented as a \textbf{directed acyclic graph} (DAG)
\begin{itemize}
  \item An \textbf{ordered sequence of nodes} with \textbf{directed edges}
  \item Nodes hold intermediate values $\zv_i$ according to operations $h_i(\cdot)$ of edges feeding into them
\end{itemize}
%TODO-BB: die "nodes" sollen hier die layers sein? dann klar sagen, denn sonst denkt man "neurons"
\item Layer $i$ maps $\zv_{i-1}$ to $\zv_i = h_i(\zv_{i-1}; \thetav_i)$
\begin{itemize}
  \item e.g. $\zv_i = \sigma(\bm{W}_i \zv_{i-1} + \bm{b}_i)$ with $\thetav_i = (\bm{W}_i, \bm{b}_i)$
\end{itemize}
\item A DNN can be represented as a \textbf{directed acyclic graph} (DAG)
\begin{itemize}
  \item An \textbf{ordered sequence of nodes} with \textbf{directed edges}
  \item Nodes hold intermediate values $\zv_i$ according to operations $h_i(\cdot)$ of edges feeding into them
\end{itemize}
\item Layer $i$ maps $\zv_{i-1}$ to $\zv_i = h_i(\zv_{i-1}; \thetav_i)$
\begin{itemize}
  \item e.g. $\zv_i = \sigma(\bm{W}_i \zv_{i-1} + \bm{b}_i)$ with $\thetav_i = (\bm{W}_i, \bm{b}_i)$
\end{itemize}
\item A DNN can be represented as a \textbf{directed acyclic graph} (DAG): an ordered sequence of \textbf{nodes} with \textbf{directed edges}
\item Each \textbf{node} $\zv_i \in \R^{p_i}$, $i = 0, \ldots, m$, collects the intermediate values of all $p_i$ neurons of the $i$-th layer \\ {\footnotesize(The input acts as the first node: $\zv_0 = \xv$)}
\item The \textbf{directed edge} into node $i$ is the $i$-th layer: a \textbf{function} mapping the previous node to the next,
$$
h_i(\cdot \,; \thetav_i): \R^{p_{i-1}} \to \R^{p_i}, \qquad \zv_i = h_i(\zv_{i-1}; \thetav_i)
$$
{\footnotesize($\thetav_i \in \R^{d_i}$ is the parameter vector of layer $i$; stacking all layers yields $\thetav = (\thetav_1, \ldots, \thetav_m)$ with $d = \sum_{i=1}^{m} d_i$)}
\end{itemize}
\vfill
\begin{center}
\begin{tikzpicture}[>=stealth', node distance=1.2cm, every node/.style={font=\footnotesize}]
  \tikzstyle{nd}=[circle, draw, very thick, minimum size=0.75cm, fill=black!4]
  \node[nd] (x) {$\xv$};
  \node[nd, right=of x] (z1) {$\zv_1$};
  \node[nd, right=of z1] (z2) {$\zv_2$};
  \node[right=0.8cm of z2] (dots) {$\cdots$};
  \node[nd, right=0.9cm of dots] (zm) {$\zv_m$};
  \node[nd, right=of zm, fill=blue!10] (L) {$\risket$};
  \draw[->,thick] (x)--(z1) node[midway,above,font=\scriptsize]{$h_1$};
  \draw[->,thick](z1)--(z2) node[midway,above,font=\scriptsize]{$h_2$};
  \draw[->,thick](z2)--(dots);
  \draw[->,thick](dots)--(zm) node[midway,above,font=\scriptsize]{$h_m$};
  \draw[->,thick](zm)--(L);
  \node[font=\scriptsize, below=0.25cm of z1] (t1){$\thetav_1$}; \draw[->,gray](t1)--(z1);
  \node[font=\scriptsize, below=0.25cm of z2] (t2){$\thetav_2$}; \draw[->,gray](t2)--(z2);
  \node[font=\scriptsize, below=0.25cm of zm] (tm){$\thetav_m$}; \draw[->,gray](tm)--(zm);
  \node[font=\scriptsize, above=0.15cm of zm]{$= f(\xv \mid \thetav)$};
  \node[font=\scriptsize, above=0.15cm of L]{$L(y, \cdot)$};
\end{tikzpicture}
\end{center}
\end{framev}


\begin{framei}{Gradient via the chain rule}
\item The risk gradient w.r.t. the parameter vector $\thetav_i$ of layer $i$ factorizes via the \textbf{chain rule}:
\begin{align*}
\nabla_{\thetav_i} \risket &= \pd{\risket}{\zv_m} \pd{\zv_m}{\zv_{m-1}} \cdots \pd{\zv_{i+1}}{\zv_i} \; \pd{\zv_i}{\thetav_i} \\
&= \underbrace{\pd{\risket}{\zv_i}}_{=: \, \adj{i} \, \in \, \R^{1 \times p_i}} \; \underbrace{\pd{\zv_i}{\thetav_i}}_{\in \, \R^{p_i \times d_i}} \; \in \R^{1 \times d_i}
\end{align*}
\item[$\Rightarrow$] A complicated gradient becomes a \textbf{product of simple, local derivatives}
\item[$\Rightarrow$] Each layer's gradient $\nabla_{\thetav_i} L$ requires its associated \textbf{`node adjoint'} $\adj{i} \in \R^{1 \times p_i}$
% TODO-BB: sollte z-bar hier wirklich eine matrix zeilenvek sein...? nochmal checken
\item[$\Rightarrow$] Each layer's gradient $\nabla_{\thetav_i} L$ requires its associated \textbf{`node adjoint'} $\adj{i} \in \R^{1 \times p_i}$
\item[$\Rightarrow$] Each layer's risk gradient $\nabla_{\thetav_i} \risket$ requires its associated \textbf{`node adjoint'} $\adj{i}$
\end{framei}

\begin{framev}{Adjoints as sequences of VJPs}
\begin{itemize}
% wir sollten diese sachen sagen: (vermutlich ganz am anfang)  
% - wir können jeden algebraischen ausdruck so darstellen
% - wir können immer die layer in jedem DAG identifizieren
% - wir betrachten jetzt jedes layer als eine funktion g_M: R^M --> R^K
% K und M sind die anzahlen der nodes pro layer, sollten wir auch notation für haben
% auf dieser folie leiten wir nämlich dann diese funktion ab....
% also "Jacobian" ist dann "Jacobian von g_M"
% ich glaub man muss hier auch aufpassen ob man über "ableitungen als funktion" redet
% bzw "ausgewertete ableitungen an einem punkt"....!
\item All node adjoints $\adj{i}$ can be recursively computed in reverse order for $i = m, (m-1), \ldots, 1$:
\item Starting from $\adj{m} = \pd{\risket}{\zv_m}$, all remaining node adjoints $\adj{i}$ can be recursively computed in reverse order for $i = m-1, m-2, \ldots, 1$:
$$
\adj{i} = \pd{\risket}{\zv_{i+1}} \pd{\zv_{i+1}}{\zv_i} = \adj{i+1} \, \pd{\zv_{i+1}}{\zv_i}
$$
\begin{itemize}
  \item $\adj{i+1} \in \R^{1 \times p_{i+1}}$ is a (row) \textbf{vector}
  \item $\pd{\zv_{i+1}}{\zv_i} \in \R^{p_{i+1} \times p_i}$ is a \textbf{Jacobian}
\end{itemize}
\item[$\Rightarrow$] This constitutes a sequence of \textbf{vector-Jacobian products (VJPs)}
\item[$\Rightarrow$] Each $\adj{i}$ must only be computed \textbf{once} and can be \textbf{reused} for lower layers
\end{itemize}
\end{framev}


\begin{framei}{Backpropagation algorithm}
\item This motivates the \textbf{backpropagation} algorithm to compute the full gradient $\nabla_{\thetav} \risket$ in \textbf{two phases}:
%
\item \textbf{Forward pass}
\begin{itemize}
  \item Set $\zv_0 = \xv$ and walk forward through the graph up to $\risket$ %via $\zv_i = h_i(\zv_{i-1}; \thetav_i)$,
  \item While doing so, \textbf{cache} every node value $\zv_i$, because the derivatives in the subsequent phase are evaluated at these values
\end{itemize}
\item \textbf{Backward pass}
\begin{itemize}
\item Compute the adjoints $\adj{i}$ as a sequence of VJPs by moving backward through the graph
\item Obtain the risk gradient w.r.t. each layer's parameter vector $\thetav_i$ as $\nabla_{\thetav_i} \risket = \adj{i} \, \pd{\zv_i}{\thetav_i}$ (again a VJP)
\end{itemize}
% \spacer
\end{framei}


\begin{framei}{Automatic differentiation}
\item (Reverse-mode) automatic differentiation (\textbf{autodiff}) efficiently automates the backpropagation algorithm
\item Traces code to automatically build a computational graph
\begin{itemize}
  \item Decomposes DNNs into \textbf{primitive operations} (e.g. $+$, $\cdot$, $\sigma$, $\exp$, \ldots; more granular than entire layers) \\
  {\footnotesize(Layer $h_i: \R^{p_{i-1}} \to \R^{p_i}$ becomes $h_i = h_{i,k} \circ \cdots \circ h_{i,1}$ with $h_{i,1}: \R^{p_{i-1}} \to \R^{p_{i,1}}$, \ldots, $h_{i,k}: \R^{p_{i,k-1}} \to \R^{p_i}$)}
\end{itemize}
\item Automatically computes derivatives
\begin{itemize}
  \item Applies the chain rule on the (more granular) graph and efficiently computes VJPs in the backward pass via rules (see next slides) 
\end{itemize}
\item[$\Rightarrow$] Autodiff is a key component powering modern ML frameworks (e.g. \texttt{autograd} in PyTorch)
\end{framei}

\begin{framei}{VJP rules}
\item Consider a primitive operation $g: \R^{q_{\mathrm{in}}} \to \R^{q_{\mathrm{out}}}$ and the VJP $\bm{v} \, \pd{g(\zv)}{\zv} \in \R^{1 \times q_{\mathrm{in}}}$
\begin{itemize}
  \item $\bm{v} \in \R^{1 \times q_{\mathrm{out}}}$ as an arbitrary (row) vector
  \item $\pd{g(\zv)}{\zv}= \bm{J}_{g}(\zv) \in \R^{q_{\mathrm{out}} \times q_{\mathrm{in}}}$ as the Jacobian of the operation evaluated at $\zv$
\end{itemize}
{\footnotesize(E.g. $\bm{v} = \adj{i}$ may be the adjoint of node $\zv_i$ and $g = h_{i,k}$ the last primitive operation of layer $i$)}
\item Autodiff stores efficient computation rules for these VJPs, i.e. $\text{vjp}_g(\bm{v}, \zv, g(\zv))= \bm{v} \, \pd{g(\zv)}{\zv}$
\begin{itemize}
  \item Reformulation that depends only on $\bm{v}$, the operation input $\zv$ and (cached) output $g(\zv)$
  \item[$\Rightarrow$] The (potentially large) Jacobian $\bm{J}_{g}(\zv)$ need not be materialized and is only implicitly multiplied
\end{itemize}
\end{framei}

\begin{framei}{VJP rules: example}
\item \textbf{Example}: Consider the \textbf{elementwise sigmoid} $g = \sigma: \R^{q} \to \R^{q}$
\begin{itemize}
  \item i.e. $\bm{a} := \sigma(\zv)$ with $a_j = \sigma(z_j)$, each output depending only on its own input
\end{itemize}
\item The Jacobian $\bm{J}_{\sigma}(\zv) = \pd{\sigma(\zv)}{\zv} \in \R^{q \times q}$ is \textbf{diagonal}
\begin{itemize}
\item $[\bm{J}_{\sigma}(\zv)]_{jj} = \sigma'(z_j) = a_j(1 - a_j)$ and $[\bm{J}_{\sigma}(\zv)]_{jk} = 0$ for $j \neq k$
\end{itemize}
%
\item Thus the VJP rule collapses to an elementwise product:
\begin{itemize}
\item $\text{vjp}_\sigma(\bm{v}, \zv, \bm{a}) = \bm{v} \, \pd{\sigma(\zv)}{\zv} = \bm{v} \odot \bm{a}^\top \odot (1 - \bm{a})^\top$
\end{itemize}
\item[$\Rightarrow$] $\mathcal{O}(q^2)$ cost for building, storing and multiplying the explicit Jacobian drops to $\mathcal{O}(q)$
\end{framei}

\begin{framev}{Worked example: forward pass}
\begin{itemize}
\item Consider a two-layer DNN with $\thetav = (\thetav_1, \thetav_2)$, $\thetav_i = (\bm{W}_i, \bm{b}_i)$, $\bm{W}_i \in \R^{p_i \times p_{i-1}}$, $\bm{b}_i \in \R^{p_i}$, scalar output ($p_2 = 1$) and squared loss $\risket = \tfrac12 (\zv_2 - y)^2$ for one observation $(\xv, y)$
\item Each layer consists of two primitive operations, $h_i = h_{i,2} \circ h_{i,1}$:
\begin{itemize}
  \item Affine map $h_{i,1}(\zv_{i-1}) = \bm{W}_i \zv_{i-1} + \bm{b}_i$ with intermediate value $\zv_{i,1} \in \R^{p_i}$
  \item Elementwise sigmoid $h_{i,2} = \sigma$
\end{itemize}
\end{itemize}
\vfill
\begin{center}
\begin{tikzpicture}[>=stealth', node distance=1.05cm, every node/.style={font=\footnotesize}]
  \tikzstyle{nd}=[circle, draw, very thick, minimum size=0.75cm, fill=black!4]
  \tikzstyle{im}=[circle, draw, thick, dashed, minimum size=0.75cm, fill=black!2]
  \node[nd] (x) {$\xv$};
  \node[im, right=of x] (z11) {$\zv_{1,1}$};
  \node[nd, right=of z11] (z1) {$\zv_1$};
  \node[im, right=of z1] (z21) {$\zv_{2,1}$};
  \node[nd, right=of z21] (z2) {$\zv_2$};
  \node[nd, right=of z2, fill=blue!10] (L) {$\risket$};
  \draw[->,thick] (x)--(z11) node[midway,above,font=\scriptsize]{$h_{1,1}$};
  \draw[->,thick] (z11)--(z1) node[midway,above,font=\scriptsize]{$h_{1,2}$};
  \draw[->,thick] (z1)--(z21) node[midway,above,font=\scriptsize]{$h_{2,1}$};
  \draw[->,thick] (z21)--(z2) node[midway,above,font=\scriptsize]{$h_{2,2}$};
  \draw[->,thick] (z2)--(L);
  \node[font=\scriptsize, below=0.25cm of z11] (t1){$\thetav_1$}; \draw[->,gray](t1)--(z11);
  \node[font=\scriptsize, below=0.25cm of z21] (t2){$\thetav_2$}; \draw[->,gray](t2)--(z21);
\end{tikzpicture}
\end{center}
\vfill
\end{framev}


\begin{framev}{Worked example: backward pass}
\begin{itemize}
\item \textbf{Adjoints} via VJPs in reverse order:
\vspace{-0.6em}
\begin{align*}
\adj{2} &= \pd{\risket}{\zv_2} = \zv_2 - y && \in \R \\
\adj{2,1} &= \adj{2} \, \pd{\zv_2}{\zv_{2,1}} = \adj{2} \, \zv_2 (1 - \zv_2) && \in \R \\
\adj{1} &= \adj{2,1} \, \pd{\zv_{2,1}}{\zv_1} = \adj{2,1} \, \bm{W}_2 && \in \R^{1 \times p_1} \\
\adj{1,1} &= \adj{1} \, \pd{\zv_1}{\zv_{1,1}} = \adj{1} \odot \zv_1^\top \odot (1 - \zv_1)^\top && \in \R^{1 \times p_1}
\end{align*}
\item \textbf{Risk gradients} w.r.t. $\thetav_i = (\bm{W}_i, \bm{b}_i)$, again via VJPs (with $\zv_0 = \xv$):
\vspace{-0.6em}
\begin{align*}
\nabla_{\bm{b}_i} \risket &= \adj{i,1} \, \pd{\zv_{i,1}}{\bm{b}_i} = \adj{i,1} && \in \R^{1 \times p_i} \\
\nabla_{\bm{W}_i} \risket &= \adj{i,1} \, \pd{\zv_{i,1}}{\bm{W}_i} = \adj{i,1}^\top \, \zv_{i-1}^\top && \in \R^{p_i \times p_{i-1}} \; \text{\footnotesize(shape of $\bm{W}_i$)}
\end{align*}
\end{itemize}
\end{framev}


% TODO-BB: in JEDEM FALL! brauchen wir hier ein längeres bsp, gezeichnet, mit mehr als einem layer...

\endlecture
\end{document}
